\documentclass[tikz,border=0pt]{standalone}
\usepackage{fontspec}
\usetikzlibrary{arrows.meta}

\setmainfont{Arial}
\setsansfont{Arial}

\definecolor{pagebg}{HTML}{F8FAFC}
\definecolor{border}{HTML}{CBD5E1}
\definecolor{textmain}{HTML}{0F172A}
\definecolor{textmuted}{HTML}{475569}
\definecolor{blue}{HTML}{2563EB}
\definecolor{bluefill}{HTML}{DBEAFE}
\definecolor{cyanfill}{HTML}{E0F2FE}
\definecolor{red}{HTML}{DC2626}
\definecolor{redfill}{HTML}{FEE2E2}
\definecolor{green}{HTML}{15803D}
\definecolor{greenfill}{HTML}{DCFCE7}

\begin{document}
\begin{tikzpicture}[x=1mm,y=1mm,line cap=round,line join=round,font=\sffamily]
  \fill[pagebg] (0,0) rectangle (360,230);
  \draw[fill=white,draw=border,line width=0.55mm,rounded corners=4mm]
    (6,6) rectangle (354,224);

  \node[font=\sffamily\bfseries\fontsize{18}{22}\selectfont,text=textmain]
    at (180,208) {Как передаётся внешнее давление};
  \node[font=\sffamily\fontsize{10.5}{13}\selectfont,text=textmuted]
    at (180,192) {Поршень одинаково увеличивает давление во всём объёме жидкости или газа};

  % Сосуд с поршнем
  \draw[fill=pagebg,draw=border,line width=0.45mm,rounded corners=2.5mm]
    (14,45) rectangle (228,180);

  \draw[fill=cyanfill,draw=textmuted,line width=0.75mm,rounded corners=2mm]
    (30,61) rectangle (212,141);
  \draw[fill=white,draw=white,line width=0.8mm]
    (91,140) rectangle (151,147);
  \draw[fill=border,draw=textmuted,line width=0.7mm]
    (94,141) rectangle (148,153);
  \draw[draw=textmuted,line width=1.0mm] (121,153) -- (121,169);

  \draw[-{Stealth[length=3.6mm,width=2.6mm]},draw=red,line width=1.1mm]
    (121,178) -- (121,160);
  \node[anchor=east,font=\sffamily\bfseries\fontsize{9.2}{11.3}\selectfont,text=red]
    at (115,170) {внешнее действие};
  \node[anchor=west,font=\sffamily\fontsize{8.8}{10.8}\selectfont,text=textmuted]
    at (153,147) {поршень};

  \node[font=\sffamily\bfseries\fontsize{11}{13.5}\selectfont,text=blue]
    at (121,130) {жидкость или газ};

  % Три произвольные точки среды
  \foreach \x/\y/\label in {68/108/A,166/111/B,119/82/C} {
    \node[circle,fill=white,draw=blue,line width=0.7mm,minimum size=12mm,inner sep=0pt,
      font=\sffamily\bfseries\fontsize{10.5}{12.5}\selectfont,text=blue]
      at (\x,\y) {\label};
  }
  \node[font=\sffamily\fontsize{8.8}{10.8}\selectfont,text=textmuted]
    at (121,69) {в каждой точке — одна и та же добавка давления};

  % Силы давления на стенки: равные короткие стрелки
  \draw[-{Stealth[length=3.2mm,width=2.3mm]},draw=blue,line width=0.95mm]
    (30,96) -- (18,96);
  \draw[-{Stealth[length=3.2mm,width=2.3mm]},draw=blue,line width=0.95mm]
    (212,96) -- (224,96);
  \draw[-{Stealth[length=3.2mm,width=2.3mm]},draw=blue,line width=0.95mm]
    (82,61) -- (82,49);
  \draw[-{Stealth[length=3.2mm,width=2.3mm]},draw=blue,line width=0.95mm]
    (160,61) -- (160,49);

  % Смысл закона
  \draw[fill=white,draw=border,line width=0.45mm,rounded corners=2.5mm]
    (238,45) rectangle (346,180);

  \draw[fill=bluefill,draw=blue!35,line width=0.4mm,rounded corners=1.8mm]
    (246,141) rectangle (338,173);
  \node[font=\sffamily\bfseries\fontsize{10.5}{13}\selectfont,text=blue]
    at (292,164) {Во все точки};
  \node[align=center,text width=84mm,font=\sffamily\fontsize{8.5}{10.5}\selectfont,text=textmain]
    at (292,150) {A, B и C получают одинаковое добавочное давление};

  \draw[fill=greenfill,draw=green!35,line width=0.4mm,rounded corners=1.8mm]
    (246,99) rectangle (338,131);
  \node[font=\sffamily\bfseries\fontsize{10.5}{13}\selectfont,text=green]
    at (292,122) {По всем направлениям};
  \node[align=center,text width=84mm,font=\sffamily\fontsize{8.5}{10.5}\selectfont,text=textmain]
    at (292,108) {Среда действует на площадки любой ориентации};

  \draw[fill=redfill,draw=red!30,line width=0.4mm,rounded corners=1.8mm]
    (246,57) rectangle (338,89);
  \node[font=\sffamily\bfseries\fontsize{10.5}{13}\selectfont,text=red]
    at (292,80) {Без ослабления};
  \node[align=center,text width=84mm,font=\sffamily\fontsize{8.5}{10.5}\selectfont,text=textmain]
    at (292,66) {Добавка не уменьшается с расстоянием от поршня};

  \draw[fill=bluefill,draw=blue!45,line width=0.45mm,rounded corners=2mm]
    (14,14) rectangle (346,36);
  \node[font=\sffamily\bfseries\fontsize{9.6}{11.8}\selectfont,text=blue]
    at (180,25) {Давление — скаляр; стрелки показывают силы давления на стенки};
\end{tikzpicture}
\end{document}
